% Physical pages 3 and 4 of the concise study: the movement of the formulary,
% stated as an argument and not as a summary of what follows. The developed
% comparison of the three interpretations begins on the fifth page.

\section{The Propers: Themes and Movement}
\zlabel{triptych:concise:themes:start}

This Mass answers a trap about one debt with a sentence about two, and it
surrounds that sentence with texts about what is owed, to whom, under whose
sovereignty, and from what kind of heart. It begins with a man who calls on
God because God hears him and ends by asking to be taught by eternal things;
between the two it sets a pagan king whose hand the Lord grasps, nations
summoned to bring the Lord glory, a church whose gospel came in power and not
in word only, and, at the centre, men who praise Jesus truthfully in order to
catch him.

The Entrance Antiphon is drawn from Psalm 17 (16), the prayer of a man who
asks God to hear a just cause. It begins where the speaker already has an
answer, \latin{Ego clamavi, quoniam exaudisti me, Deus}: I have called because
thou hast heard. He asks God to bend his ear, to keep him ``as the apple of
thy eye'', and to protect him ``under the shadow of thy wings'' (Ps 16:6, 8;
psalm verses quoted from the Douay--Rheims bear their Vulgate numbers).
The psalm around these lines, unsung, prays from lips without deceit and from
a heart God has proved; the same v.~8 opens by asking protection ``From them
that resist thy right hand'', words the antiphon leaves out. The Collect then
makes two requests of the almighty and eternal God: that the assembly always,
\latin{semper}, hold its will devoted to him, and that it serve his majesty,
\latin{maiestati tuae}, \latin{sincero corde}, with an undivided heart. Will,
service and heart are named before any reading is heard.

The first reading is an oracle addressed to a foreigner. ``Thus saith the Lord
to my anointed Cyrus, whose right hand I have taken hold of, to subdue nations
before his face'' (Is 45:1). The Lectionary leaves out the promises of broken
gates and hidden treasures and resumes with the reason: for the sake of Jacob
his servant the Lord has called Cyrus by name, ``and thou hast not known me'';
there is no God besides the Lord, who girded the king though the king did not
know him, so that from the rising of the sun to the west all may know ``that
there is none besides me'' (45:4--6). The title Israel gave its kings is given
to a Persian who worships other gods, and the Lord's purpose is not the
Persian's. The \work{Ordo lectionum Missae} sets a title over each reading and
explains that in Ordinary Time the Old Testament reading is chosen for its
relation to the Gospel (\work{Praenotanda} 105--107). Over this one it writes
\latin{Apprehendi dexteram Cyri ut subiciam ante faciem eius gentes}; over the
Gospel, \latin{Reddite quae sunt Caesaris Caesari, et quae sunt Dei Deo}. The
two titles that belong together each name a gentile ruler.

The psalm responds by turning ``there is none besides me'' into a summons to
the nations. Sing to the Lord a new song, all the earth; declare his glory
among the gentiles; he is to be feared above all gods, ``For all the gods of
the Gentiles are devils: but the Lord made the heavens'' (Ps 95:5), where the
Hebrew calls them idols. The kindreds of the peoples are to bring the Lord
glory and to come into his courts with gifts, and to say among the nations,
``the Lord hath reigned''; the strophe ends with the promise that ``he will
judge the people with justice'' (95:10).

The second reading opens the Lectionary's course through the First Letter to
the Thessalonians, a course independent of the Gospel's. Paul, with Silvanus
and Timothy, writes ``to the church of the Thessalonians: in God the Father and
in the Lord Jesus Christ'', remembers its work of faith, labour of charity and
enduring hope, and knows its election because ``our gospel hath not been unto
you in word only, but in power also: and in the Holy Ghost and in much
fulness'' (1~Thes 1:1--5). Acts, which is not read, tells how that church was
founded amid a tumult in which believers were dragged before the city's rulers
on the charge of saying ``that there is another king, Jesus'' (Acts 17:7). The
Alleluia verse greets the Gospel with Paul's words to the Philippians, that
they ``shine as lights in the world'', holding forth the word of life (Phil
2:15--16). The words stand at the end of a sentence that begins with the wish
that they be ``blameless and sincere children of God'' in the midst of a
crooked and perverse generation; the Latin behind ``sincere'' is
\latin{simplices}, single.

The Gospel is a trap and its undoing. In the temple at Jerusalem, after the
parable of the husbandmen, which the chief priests and Pharisees knew he spoke
of them (Mt 21:45), and the parable of the wedding feast, the Pharisees take
counsel how to ensnare Jesus in his speech, and send their disciples with the
Herodians. They open with praise: ``Master, we know that
thou art a true speaker and teachest the way of God in truth. Neither carest
thou for any man: for thou dost not regard the person of men'' (Mt 22:16).
Every clause is true. Then the question, built to wound on either answer: ``Is
it lawful to give tribute to Caesar, or not?'' Jesus, ``knowing their
wickedness'', calls them hypocrites and asks for ``the coin of the tribute''.
They bring a denarius. He asks whose image and inscription it bears, and they
answer, ``Caesar's''. Then: ``Render therefore to Caesar the things that are
Caesar's; and to God, the things that are God's.'' They had asked whether it
was lawful to give, \latin{dare}; he commands them to render,
\latin{reddite}, to give back what is owed. And the second clause, about which
nobody asked, names no object at all. The reading stops there, before the
questioners wonder and go away.

The Prayer over the Offerings calls the gifts on the altar God's own,
\latin{donis tuis}, before they are offered, and asks that those who serve them
may do so \latin{libera mente}, with a free mind, and be cleansed by the very
mysteries in which they minister. At Communion the Missal offers two
antiphons. The first takes from Psalm 33 (32) the eyes of the Lord that rest on
those who fear him and hope in his mercy; just before it the psalm says, in
verses not sung, that the Lord looks on all the sons of men, ``He who hath made
the hearts of every one of them'', and that ``The king is not saved by a great
army'' (Ps 32:15--16). The second keeps only the end of Mark 10:45: the Son of
Man came to give his life as a ransom for many. In Mark the sentence answers
the disciples' quarrel about the first places and follows the words ``they who
seem to rule over the Gentiles lord it over them \ldots\ But it is not so among
you'' (10:42--43), which the antiphon does not sing. The Prayer after Communion
asks that frequent sharing in heavenly things make the assembly advance, aided
by temporal benefits and instructed, \latin{erudiamur}, by eternal ones.

The sentence at the centre leaves three questions open, and each gives one
interpretation of the whole formulary, developed in the full study with four
senses of its own.

The first, the coin and the image, asks what ``the things that are God's''
are, and is heard first from Jesus' question about the image. St Augustine
answers that as the coin bears Caesar's image, so man bears God's, and must
render himself: ``to Caesar his money, to God yourselves''. St Ambrose, on
Luke's form of the saying, adds that Christ bears no image of Caesar because
he is the image of God. St Hilary of Poitiers reaches the debt of the whole
self from another premise, that body, soul and will are owed to the God from
whom they come; St Thomas Aquinas keeps a literal sense and gives the image
reading as the mystical one. St Jerome reads the second clause of the saying
literally, of tithes and offerings.

The second, Cyrus and Caesar, asks how God's people can owe anything to a
ruler who does not know God, and how far the debt runs. It is heard from the
Gospel's command with the oracle the \work{Ordo} sets beside it. Jerome and
Aquinas explain that Cyrus is called the Lord's anointed for his royal
dignity; St John Chrysostom says that tribute is rendered as a debt and ends
where godliness would be harmed; Augustine tells of Christian soldiers who
served an unbelieving emperor and refused him only in the cause of Christ.
Hilary and Ambrose make the debt follow the possession of the world's goods,
while Chrysostom holds even monks and Apostles bound.

The third, the sincere heart, asks why men who speak the truth are called
hypocrites, and is heard from Jesus' rebuke. Jerome, Chrysostom, Aquinas and
Augustine each place the fault in the distance between the questioners'
words and their hearts. Chrysostom, Hilary and Aquinas make Christ's knowledge
of that heart a sign of who he is; Jerome states the same discernment as the
first virtue of anyone who answers.

The three hold much in common. All read the saying as setting Caesar's claim
beneath God's, and none as a division of life into two equal realms. All rest
on the whole scene in the temple, from the plot to the verdict, and all hear
the Collect's will, majesty and heart, each for a reason of its own. They
differ in the words each hears first, in what each needs from the rest of the
Mass, and in the difficulty each must meet; and at two points, what God is
owed and what the debt to Caesar rests on, the witnesses themselves differ. The
commentary that follows sets the interpretations' answers beside one another
where the texts put the same question to all three.
\zlabel{triptych:concise:themes:end}
